% Generated by code/tables/tabA-measurement.py. Do not edit by hand.

\begin{table}[htbp]
\centering
\begin{threeparttable}
\caption{Capability-only and constraint-adjusted scores across tasks}
\label{tab:techfeas}
\footnotesize
\begin{spacing}{1}
\begin{tabular}{@{}lcccc@{}}
\toprule
 \\
\midrule
\multicolumn{5}{@{}l}{\emph{A. Three-model score, by stage}} \\
 & Capability only & Constraint adjusted & Pearson & Spearman \\
Stage 1 & 16.1 & 11.0 & 0.975 & 0.962 \\
Stage 2 & 36.1 & 25.5 & 0.964 & 0.965 \\
Stage 3 & 46.6 & 35.0 & 0.961 & 0.962 \\
Stage 4 & 55.4 & 43.6 & 0.969 & 0.972 \\
Change, stage 1 to 4 & 39.3 & 32.5 & 0.933 & 0.928 \\
\addlinespace
\multicolumn{5}{@{}l}{\emph{B. Each model, stage 4}} \\
Claude Haiku 4.5 & 59.8 & 47.1 & 0.928 & 0.930 \\
GPT-5.6-luna & 53.9 & 40.9 & 0.948 & 0.951 \\
Gemini 3.1 Flash-Lite & 52.6 & 42.7 & 0.971 & 0.973 \\
\bottomrule
\end{tabular}
\end{spacing}
\begin{tablenotes}\notesize
\item Notes: Ratings A and B were each given in two versions: the capability of the technology alone, and the rating after the economic, regulatory, and organizational constraints on using AI for the subtask, which the paper uses. Rating C has one version and enters both composites. Means are across the 19,278 tasks with all three models; correlations are across tasks. The change rows compare each version's change from stage 1 to stage 4.
\end{tablenotes}
\end{threeparttable}
\end{table}
